%%%%%%%%%%%%%%%%%%%%%%%%%%%% Document Setup %%%%%%%%%%%%%%%%%%
\documentclass[10pt]{article}
\usepackage[T2A,T1]{fontenc}
\usepackage[utf8]{inputenc}
\usepackage[russian,english]{babel}
\usepackage{CJKutf8}
\AtBeginDvi{\input{zhwinfonts}}
%\usepackage{polyglossia}
\makeatletter
\newlength{\bibhang}
\setlength{\bibhang}{1em}
\newlength{\bibsep}
 {\@listi \global\bibsep\itemsep \global\advance\bibsep by\parsep}
%\newenvironment{bibsection}%
 %       {\vspace{-\baselineskip}\begin{list}{}{%
   %    \setlength{\leftmargin}{\bibhang}%
     %  \setlength{\itemindent}{-\leftmargin}%
       %\setlength{\itemsep}{\bibsep}%
       %\setlength{\parsep}{\z@}%
        %\setlength{\partopsep}{0pt}%
        %\setlength{\topsep}{0pt}}}
        %{\end{list}\vspace{-.6\baselineskip}}
\makeatother
\reversemarginpar
\usepackage{geometry}\geometry{a4paper,total={170mm,257mm},left=1.18in,right=1.18in,top=0.79in,bottom=0.79in}
%\usepackage[paper=letterpaper,
           %\includefoot, % Uncomment to put page number above margin
   %         marginparwidth=0.6in, % Length of section titles
         %   marginparsep=0.05in, % Space between titles and text
      %      margin=1.18in, % 1 inch margins
            %includemp]{geometry}
%\usepackage{eurosans}
\setlength{\parindent}{0in}
\usepackage{enumerate} 
\usepackage{etaremune}
\usepackage{calc}
\usepackage{bbding}
\usepackage{paralist}
\usepackage{fancyhdr,lastpage}
\pagestyle{fancy}
%\pagestyle{empty} % Uncomment this to get rid of page numbers
\fancyhf{}\renewcommand{\headrulewidth}{0pt}
\fancyfootoffset{\marginparsep+\marginparwidth}
\newlength{\footpageshift}
\setlength{\footpageshift}
          {0.5\textwidth+1.\marginparsep+0.5\marginparwidth-2in}
\lfoot{\hspace{\footpageshift}%
       \parbox{4in}{\, \hfill %
                    \arabic{page} of \protect\pageref*{LastPage} % +LP
% \arabic{page} % -LP
                    \hfill \,}}
\usepackage{color,hyperref}
\usepackage{graphics}
\usepackage{wrapfig}
\usepackage{longtable}
\usepackage[usenames,dvipsnames]{xcolor}
\definecolor{darkblue}{rgb}{0.0,0.0,0.3}
\definecolor{mygreen}{RGB}{44,204,148}
\definecolor{darkgreen}{RGB}{13,107,27}
\definecolor{lightgreen}{RGB}{0,170,37}
\definecolor{provinces}{RGB}{11,100,132}
\definecolor{mypurple}{RGB}{77,24,137}
\definecolor{haiyan}{RGB}{8,150,178}
\definecolor{anticred}{RGB}{165,62,11}
\hypersetup{colorlinks,breaklinks,
            linkcolor=darkblue,urlcolor=darkblue,
            anchorcolor=darkblue,citecolor=darkblue}
%%%%%%%%%%%%%%%%%%%%%%%% End Document Setup %%%%%%%%%%%%%%%%%%%
%%%%%%%%%%%%%%%%%%%%%%%%%%% Helper Commands %%%%%%%%%

\newcommand{\makeheading}[1]%
        {\hspace*{-\marginparsep minus \marginparwidth}%
         \begin{minipage}[t]{\textwidth+\marginparwidth+\marginparsep}%
                {\large \bfseries #1}\\[-0.15\baselineskip]%
                 \rule{\columnwidth}{1pt}%
         \end{minipage}}
\renewcommand{\section}[2]%
        {\pagebreak[2]\vspace{1.3\baselineskip}%
         \phantomsection\addcontentsline{toc}{section}{#1}%
         \hspace{0in}%
         \marginpar{
         \raggedright \scshape #1}#2}
% An itemize-style list with lots of space between items
\newenvironment{outerlist}[1][\enskip\textbullet]%
        {\begin{itemize}[#1]}{\end{itemize}%
         \vspace{-.6\baselineskip}}
% An environment IDENTICAL to outerlist that has better pre-list spacing
% when used as the first thing in a \section
\newenvironment{lonelist}[1][\enskip\textbullet]%
        {\vspace{-\baselineskip}\begin{list}{#1}{%
        \setlength{\partopsep}{0pt}%
        \setlength{\topsep}{0pt}}}
        {\end{list}\vspace{-.6\baselineskip}}

% An itemize-style list with little space between items
\newenvironment{innerlist}[1][\enskip\textbullet]%
        {\begin{compactitem}[#1]}{\end{compactitem}}
% An environment IDENTICAL to innerlist that has better pre-list spacing
% when used as the first thing in a \section
\newenvironment{loneinnerlist}[1][\enskip\textbullet]%
        {\vspace{-\baselineskip}\begin{compactitem}[#1]}
        {\end{compactitem}\vspace{-.6\baselineskip}}
% To add some paragraph space between lines.
% This also tells LaTeX to preferably break a page on one of these gaps
% if there is a needed pagebreak nearby.
\newcommand{\blankline}{\quad\pagebreak[2]}
% Uses hyperref to link DOI
\newcommand\doilink[1]{\href{http://dx.doi.org/#1}{#1}}
\newcommand\doi[1]{doi:\doilink{#1}}
% For \url{SOME_URL}, links SOME_URL to the url SOME_URL
\providecommand*\url[1]{\href{#1}{#1}}
% Same as above, but pretty-prints SOME_URL in teletype fixed-width font
\renewcommand*\url[1]{\href{#1}{\texttt{#1}}}

\include{tikz-3dplot}
\usepackage{tikz,tikz-3dplot}
\usepackage{amssymb}
\usepackage{xifthen}
\usepackage{pstricks,pst-solides3d}
\usetikzlibrary{chains}
\usepackage{pgfplots}
\pgfplotsset{
compat=1.16,
    compat/path replacement=1.5.1,
}
\usepackage{pgfmath}
\usetikzlibrary{matrix}
\usetikzlibrary{arrows,decorations.pathmorphing,backgrounds,positioning,fit,petri}
\usepgfmodule{decorations}
\usepgflibrary{decorations.shapes} 
\usetikzlibrary[decorations.shapes] 
\usetikzlibrary{mindmap}
\usepgflibrary{patterns}
\usetikzlibrary[patterns]
\usetikzlibrary{petri}
\usepgflibrary{plothandlers}
\usetikzlibrary{plothandlers}
\usepackage{pgfpages}
\usepgflibrary{arrows}
\usetikzlibrary{arrows}
\usetikzlibrary{trees}
\usetikzlibrary{topaths}
\usetikzlibrary{3d}
\usetikzlibrary{positioning}
\usetikzlibrary{fit}
\usepgflibrary{shapes.misc}
\usetikzlibrary{shapes.misc}
\usepgflibrary{shadings}
\usetikzlibrary{shadings}
\usepgfmodule{shapes}
\usepgfplotslibrary{external}
\usetikzlibrary{pgfplots.external}
\usepgfplotslibrary{colorbrewer}
\usepgfplotslibrary{patchplots}
\usetikzlibrary{plotmarks}
\usepackage{color} 
\usepackage{xcolor} 
%\pgfpagesuselayout{resize to}[a4paper,border shrink=20mm]
\pgfplotsset{colormap={cool}{rgb255(0cm)=(255,255,255); rgb255(1cm)=(0,128,255); rgb255(2cm)=(255,0,255)}}
  \pgfplotsset{colormap={violet}{rgb255=(25,25,122) color=(white) rgb255=(238,140,238)}}
  \pgfplotsset{  colormap={bluered}{rgb255(0cm)=(0,0,180); rgb255(1cm)=(0,255,255); rgb255(2cm)=(100,255,0);rgb255(3cm)=(255,255,0); rgb255(4cm)=(255,0,0); rgb255(5cm)=(128,0,0)}}
 \pgfplotsset{/pgfplots/colormap={winter}{rgb255=(0,0,255) rgb255=(0,255,128)}}
 \pgfplotsset{/pgfplots/colormap={spring}{rgb255=(255,0,255) rgb255=(255,255,0)}}
 \pgfplotsset{colormap={CM}{rgb=(0,0,1) color=(black) rgb255=(238,140,238)}}
 
\usetikzlibrary{spy}
 \pgfplotsset{every patch/.style={miter limit=50},}
\usepgfplotslibrary{colormaps}


\newcommand*\subfigheight{\linewidth}
\newcommand*\subfigwidth{\linewidth}

\begin{document}
\selectlanguage{russian}

\begin{tikzpicture} [spy using outlines={red,circle,magnification=1.5, size=3cm},connect spies]
  \begin{axis}[colorbar sampled line,small,title=Experiment №5807,view={240}{30},xlabel={\tiny{Time, h}},
    ylabel={\tiny{Immersion depth, mm}}, zlabel={\tiny{($C_{eq}$), MPa}}, xlabel style={sloped like x axis},
    ylabel style={sloped},grid=major,minor x tick num=10,minor y tick num=10,
    colormap/PuOr,]
    
    \addplot3+ [line join=round,miter limit=50,line join=bevel][
only marks, scatter,  ]
coordinates {
(0,0.50,0.160) (30,0.85,0.075) (60,0.92,0.072) (90,1.05,0.069) (120,1.10,0.066) (150	,1.11,0.064) (180,1.13,0.062) (210,1.15,0.059) (240,1.18,0.057) (270,1.20,0.055)
(0,0.30,0.17) (1.0,0.35,0.165) (1.5,0.37,0.16) (2,0.40,0.155) (3,0.42,0.15) (4,0.45,0.14) (5,0.5,0.13) (6,0.55,0.12) (7,0.57,0.11) (8,0.60,0.10)
(0,0.370,0.22) (1,0.375,0.21) (3,0.380,0.20) (4,0.390,0.19) (5,0.395,0.18) (6,0.400,0.17) (8,0.405,0.16) (10,0.410,0.15) (13,0.415,0.155) (15,0.420,0.14) 
(0,0.320,0.170) (0.5,0.380,0.165) (1.0,0.390,0.160) (1.5,0.400,0.155) (2,0.420,0.150) (2.5,0.440,0.140) (3,0.460,0.135) (3.5,0.480,0.127) (4,0.500,0.123) (5,0.520,0.114) 
(0,0.315,0.180) (1,0.324,0.163) (3,0.332,0.162) (4,0.408,0.158) (5,0.417,0.151) (6,0.442,0.142) (8,0.458,0.141) (10,0.473	,0.138) (13,0.502,0.136) (15,0.523,0.135) 
(0,0.372,0.21) (1,0.373,0.20) (3,0.382,0.19) (4,0.393,0.18) (5,0.393,0.17) (6,0.402,0.15) (8,0.407,0.14) (10,0.411,0.12) (13,0.417,0.11) (15,0.421,0.10) 
};
	\coordinate (spypoint) at (4,0.5,0.16);
        \coordinate (magnifyglass) at (200,1.0,0.18);
\end{axis}
	\spy on (spypoint) in node at (magnifyglass);
\end{tikzpicture}
	\hskip 10pt 	
	\begin{tikzpicture} [spy using outlines={red,rectangle,magnification=1.3, size=3.3cm},connect spies]
	\begin{axis}[colorbar sampled line,small,title=Experiment №5817,xlabel={\tiny{Time, h}},
    ylabel={\tiny{Immersion depth, mm}}, zlabel={\tiny{($C_{eq}$) , MPa}}, xlabel style={sloped like x axis},
    ylabel style={sloped},colormap name=hot,grid=major,view={240}{30},minor x tick num=10,minor y tick num=10]
    \addplot3+ [mesh,scatter,
        samples=42,domain=0:1, ] coordinates {
(0,0.50,0.150) (30,0.80,0.080) (60,0.85,0.075) (90,0.96,0.063) (120,1.00,0.062) (150,1.15,0.061) (180,1.25,0.060) (210,1.52,0.059) (240,1.80,0.058)    
(0,0.61,0.120) (2,0.66,0.110) (4,0.68,0.103) (6,0.72,0.090) (7,0.76,0.087) (8,0.78,0.083) (10,0.82,0.081) (13,0.83,0.075) (15,0.84,0.070)  
(0,0.82,0.080) (1,0.84,0.078) (2,0.86,0.074) (3,0.91,0.072) (4,0.94,0.070) (5,0.99,0.066) (6,1.13,0.064) (7,1.20,0.062) (8,1.32,0.060)  
(0,0.42,0.150) (1,0.44,0.147) (2,0.47,0.143) (3,0.52,0.132) (4,0.55,0.121) (5,0.58,0.116) (6,0.60,0.105) (7,0.62,0.099) (8,0.64,0.092)  
(0,0.40,0.111) (2,0.47,0.105) (4,0.52,0.101) (6,0.55,0.096) (8,0.61,0.092) (10,0.64,0.087) (12,0.71,0.085) (13,0.75,0.084) (15,0.81,0.082)  
(0,0.62,0.110) (2,0.66,0.103) (4,0.68,0.098) (6,0.71,0.094) (8,0.74,0.087) (10,0.76,0.085) (12,0.77,0.083) (13,0.78,0.081) (15,0.79,0.078)     
};
	\coordinate (spypoint) at (50,0.5,0.09);
        \coordinate (magnifyglass) at (200,1.5,0.12);
\end{axis}
\spy on (spypoint) in node at (magnifyglass);
\end{tikzpicture}\\


	\hskip 10pt % insert a non-breaking space of specified width.
	\begin{tikzpicture} \begin{axis}[colorbar sampled line,small,title=Experiment №5821,xlabel={\tiny{Time, h}},
    ylabel={\tiny{Immersion depth, mm}}, zlabel={\tiny{($C_{eq}$), MPa}},xlabel style={sloped like x axis},
    ylabel style={sloped},colormap/PRGn,grid=major,view={210}{30},minor x tick num=10,minor y tick num=10]
    \addplot3+ [mesh,scatter,
        samples=42,domain=0:1, ] coordinates {
 (0,0.30,0.220) (30,0.52,0.140) (60,0.58,0.120) (90,0.60,0.110) (120,0.61,0.108) (150,0.62,0.100) (180,0.63,0.095) (210,0.64,0.090)   
 (0,0.32,0.217) (2,0.35,0.203) (3,0.41,0.196) (4,0.44,0.174) (5,0.46,0.156) (6,0.48,0.142) (7,0.53,0.136) (8	,0.57,0.125)   
 (0,0.51,0.153) (2,0.54,0.126) (3,0.57,0.112) (4,0.59,0.103) (5,0.61,0.096) (6,0.63,0.091) (7,0.64,0.089) (8,0.66,0.088)   
 (0,0.50,0.154) (3,0.55,0.125) (4,0.58,0.109) (5,0.60,0.102) (6,0.62,0.099) (8,0.64,0.097) (10,0.66,0.096) (12,0.68,0.095)   
 (0,0.51,0.178) (3,0.53,0.171) (4,0.57,0.166) (5,0.59,0.152) (6,0.61,0.143) (8,0.63,0.136) (10,0.65,0.130) (12,0.67,0.127)   
 (0,0.38,0.173) (3,0.41,0.169) (4,0.42,0.165) (5,0.44,0.151) (6,0.46,0.142) (8,0.48,0.137) (10,0.50,0.135) (12,0.52,0.134)   
};
\end{axis}

\end{tikzpicture}
	\hskip 10pt % insert a non-breaking space of specified width.
	\begin{tikzpicture} \begin{axis}[colorbar sampled line,small,title=Experiment №5828,xlabel={\tiny{Time, h}},
    ylabel={\tiny{Immersion depth, mm}}, zlabel={\tiny{($C_{eq}$) , MPa}},xlabel style={sloped like x axis},
    ylabel style={sloped},colormap/cool,grid=major,view={240}{30},minor x tick num=10,minor y tick num=10]
    \addplot3+ [mesh,scatter,
        samples=36,domain=0:1, ] coordinates {
(0,0.40,0.130) (20,0.60,0.110) (40,0.70,0.090) (60,0.90,0.080) (80,1.00,0.075) (100,1.10,0.065)    
(0,0.35,0.160) (2,0.42,0.140) (4,0.51,0.120) (6,0.53,0.115) (8,0.58,0.110) (10,0.62,0.100) 
(0,0.65,0.110) (3,0.75,0.095) (8,0.80,0.081) (14,0.85,0.070) (16,0.90,0.064) (18	,1.19,0.050) 
(0,0.55,0.150) (2,0.62,0.132) (4,0.71,0.113) (6,0.73,0.102) (8,0.78,0.090) (10,0.82,0.080) 
(0,0.41,0.162) (3,0.53,0.143) (8,0.57,0.121) (14	,0.60,0.115) (16,0.63	,0.104) (18,0.65,0.096) 
(0,1.11,0.061) (2,1.23,0.048) (4,1.28,0.041) (6,1.32,0.037) (8,1.34,0.031) (10,1.35,0.026)       
};
\end{axis}
\end{tikzpicture}
\\
%	\hskip 10pt % insert a non-breaking space of specified width.
	\begin{tikzpicture} \begin{axis}[colorbar sampled line,small,title=Experiment №5822,xlabel={\tiny{Time, h}},
    ylabel={\tiny{Immersion depth, mm}}, zlabel={\tiny{($C_{eq}$), MPa}},xlabel style={sloped like x axis},
    ylabel style={sloped},colormap/RdYlGn,grid=major,view={240}{30},minor x tick num=10,minor y tick num=10]
    \addplot3+ [mesh,scatter, samples=54,domain=0:1, ] coordinates {
(0,0.30,0.200) (30,0.61,0.100) (60,0.71,0.080) (90,0.78,0.075) (120,0.80,0.065) (150,0.82,0.050) (180,0.83,0.047) (210,0.85,0.042) (240,0.90,0.040)   
(0,0.50,0.112) (2,0.61,0.102) (4,0.68,0.084) (6,0.71,0.081) (9,0.73,0.076) (12,0.74,0.067) (13,0.75,0.064) (14,0.76,0.061) (15,0.78,0.053)   
(0,0.25,0.210) (2,0.26,0.204) (2.5,0.31,0.092) (3.2,0.34,0.083) (3.5,0.35,0.073) (4.0,0.36,0.066) (4.5,0.37,0.061) (4.8,0.38	,0.069) (5,0.41,0.056)   
(0,0.45,0.150) (2,0.46,0.137) (2.5,0.49,0.129) (3.2,0.51,0.125) (3.5,0.54,0.121) (4.0,0.56,0.115) (4.5,0.58,0.107) (4.8	,0.59	,0.102) (5,0.61,0.097)   
(0,0.41,0.144) (2,0.42,0.141) (4	,0.43,0.135) (6,0.45,0.131) (9,0.47,0.127) (12,0.50,0.123) (13,0.53,0.122) (14,0.57,0.121) (15,0.58,0.120)   
(0,0.42,0.141) (2,0.43,0.139) (4	,0.45,0.137) (6,0.47,0.133) (9,0.49,0.125) (12,0.52,0.121) (13,0.55,0.118) (14,0.59,0.116) (15,0.60,0.114)        
        };
\end{axis}
\end{tikzpicture}
	\hskip 10pt % insert a non-breaking space of specified width.
	\begin{tikzpicture} [spy using outlines={red,circle,magnification=1.5, size=3cm},connect spies]
	\begin{axis}[colorbar sampled line,small,title=Experiment №5825,xlabel={\tiny{Time, h}},
    ylabel={\tiny{Immersion depth, mm}}, zlabel={\tiny{($C_{eq}$), MPa}},xlabel style={sloped like x axis},
    ylabel style={sloped},colormap/BrBG-11,grid=major,view={240}{30},minor x tick num=10,minor y tick num=10]
    \addplot3+ [mesh,scatter, samples=54,domain=0:1, ] coordinates {
(0,0.21,0.300) (20,0.24,0.150) (40,0.27,0.135) (60,0.32,0.130) (80,0.35,0.125) (100,0.38,0.123) (120	,0.39,0.121) (140,0.42,0.120)     
(0,0.35,0.210) (2,0.37,0.185) (4,0.39,0.163) (6,0.42,0.142) (8,0.45,0.131) (10,0.47,0.119) (12,0.49,0.112) (18,0.52,0.107)   
(0,0.40,0.157) (2,0.43,0.143) (4,0.47,0.136) (6,0.51,0.128) (8,0.54,0.122) (10,0.57,0.116) (12,0.59,0.107) (18,0.61,0.103)   
(0,0.51,0.120) (2,0.62,0.112) (3,0.64,0.101) (4,0.69,0.094) (5,0.71,0.081) (6,0.75,0.074) (7,0.77,0.061) (8,0.79,0.050)   
(0,0.43,0.170) (2,0.46,0.162) (3,0.50,0.151) (4,0.54,0.144) (5,0.57,0.131) (6,0.60,0.124) (7,0.62,0.111) (8,0.64,0.105)   
(0,0.41,0.158) (2,0.44,0.144) (3,0.49,0.137) (4,0.53,0.129) (5,0.55,0.123) (6,0.58,0.117) (7,0.60,0.108) (8,0.62,0.104)   
        };
        	\coordinate (spypoint) at (4,0.6,0.16);
        \coordinate (magnifyglass) at (160,0.5,0.18);
\end{axis}
	\spy on (spypoint) in node at (magnifyglass);
\end{tikzpicture}
 

%	\begin{flushleft}
%		\selectlanguage{russian}\textbf{Рис. 2. Временные ряды изменения велиhины длительного эквивалентного сцепления грунтов ($C_{eq}$) с увелиhением глубины погружения образцов. Увелиhенный фрагмент графика: 8-hасовое эквивалентное сцепление. Серии испытаний-II (№9-14). 3D-визуализация: \LaTeX{}}\\\vspace{6pt}
%		\textbf{Fig. 2. Time series analysis of the variations of the long-term equivalent cohesion of soils ($C_{eq}$) under increasing depth of the samples. Enlarged: graphs of the 8-hours equivalent cohesion of soils. Series-II (№9-16). 3D: \LaTeX{}}
%	\end{flushleft}

\end{document}
